% ==============================================================================
%  Chapter 6 -- Validation and regression experiments
%  Purpose: the "does the architecture work at all" chapter, ~10-15 pages.
%  Contains the inductive-bias-gap motivation experiment (done, numbers
%  known) and the supervised-regression shootout GELT vs. L-CNN vs. CNN
%  (scripts exist: train_gelt.py / train_lcnn.py / train_cnn.py; the
%  systematic matched-parameter comparison still needs to be RUN).
% ==============================================================================

\section{The inductive-bias gap}
\label{sec:reg-inductive-gap}

\TODO{The motivating experiment. Setup: Haar-random $\mathbb{Z}_2$ links
in 2D, target = Wilson action. Analytic backdrop (worth spelling out):
every plaquette is $\pm 1$ with mean zero, plaquette pairs share 0 or 2
links so they are uncorrelated, hence $\mathrm{Var}(S) = L^2$ -- so
compare via $R^2 = 1 - \mathrm{MSE}/\mathrm{Var}(y)$, never raw MSE
across $L$. Result: a CNN fed PLAQUETTES reaches $R^2 \approx 0.99$
(linear sum of its inputs); the same CNN fed raw LINKS sits at
$R^2 \approx 0$. The gap is exactly what the equivariant transport
provides. Caveat paragraph: a perfect regressor on Haar data memorises
the action FUNCTION; physics ($\beta$-dependence) requires the Metropolis
data path.}

\section{Experimental protocol}
\label{sec:reg-protocol}

\TODO{Common recipe for all regression runs: ensembles (Haar vs.\
Metropolis at fixed $\beta$), train/val/test splits, target
standardisation, matched parameter counts, optimizer/scheduler, seeds.
One table.}

\section{Regression targets}
\label{sec:reg-targets}

\TODO{The three physical targets and why each probes something different:
Wilson action (linear in plaquettes -- sanity floor), rectangular Wilson
loops $R \times T$ (requires transport beyond the plaquette),
topological charge density $q(x)$ (per-site, parity-odd, the
interpretability substrate). \texttt{train\_gelt.py} currently trains on
$q(x)$.}

\section{Results: GELT vs.\ L-CNN vs.\ CNN}
\label{sec:reg-results}

\TODO{THE MATCHED-PARAMETER SHOOTOUT -- still to be run systematically.
Planned figure: $R^2$ (or standardized MSE) per target per model, plus
learning curves. Include the non-equivariant CNN as the lower anchor and
check the L-CNN reproduces Favoni-like behaviour. Also report the
\texttt{mode="average"} vs.\ \texttt{mode="single"} transport A/B here
(does path averaging dilute a specific-path target like an off-axis
Wilson loop?).}

\section{Gauge-invariance in practice}
\label{sec:reg-invariance}

\TODO{Numerical invariance drift of trained models under random and
worst-case $\Omega$ (float32 training vs.\ float64 checks) -- the
empirical companion to \cref{sec:gelt-verification}.}
